% Revision 08: reviewed initial-mass/time/event output scope; revision07 preserved.
\documentclass[11pt]{article}
\usepackage[T1]{fontenc}
\usepackage[utf8]{inputenc}
\usepackage{amsmath,amssymb,mathtools}
\usepackage{geometry}
\usepackage[hidelinks]{hyperref}
\geometry{margin=1in}
\setlength{\emergencystretch}{2em}
\title{An Energy-Based Probabilistic Framework for Fatigue Crack Initiation:\\
Governing Equations and Conservative Cell Coupling}
\author{Donghyeok Kim \quad Minsoo Kang \quad Junhyeok Park}
\date{Working manuscript --- 18 October 2026}
\begin{document}
\maketitle
\begin{abstract}
We formulate a route from a nonlinear lattice energy to the probability of
first local opening under a specified load history. The retained two-coordinate
energy drives a conservative Fokker--Planck equation; irreversible loss from
the intact population defines the first-opening cumulative probability.
We distinguish this event from experimentally detected crack initiation and
state the physical clock and observation records needed to connect them.
Exact population balances show why interacting regions cannot automatically
be replaced by independent cells. Work consistency specifies conditional
energy, volume and time mappings for coupling to spatial finite volumes.
Reviewed analytical results include a finite-time load sensitivity bound
and uniform far-field force and curvature bounds for the original lattice sum.
These results retain their explicit assumptions. Atomistic justification of
the reduced diffusion, material calibration, a complete spatial constitutive
closure and experimental crack-initiation validation remain open.
\end{abstract}
\begin{quote}\small
Draft scope: theoretical framework and conditional results. No calibrated
aluminium fatigue life or experimentally validated initiation CDF is claimed.
The earlier one-dimensional draft is preserved separately; its fitted values
and imposed spacing-distribution shapes are not current-model inputs.
\end{quote}

% Editorial integration, 2026-10-18. No new scientific calculation.
% Current-model main-body replacement; preserve the old draft in a backup.
% Title, author list, abstract and bibliography are managed by the coordinator.
\section{Introduction}
Fatigue crack initiation is a first-event problem: under a specified loading
history, what is the probability that an initially intact region has opened
by a given time or cycle? A physically grounded answer requires more than
an empirical lifetime curve. It requires an energy landscape, a law for
probability transport, an initial ensemble and an operational opening event.
Connecting that event to an observed crack additionally requires a physical
clock and a measurement model.
The Lennard--Jones family supplies a tractable interaction model, with known
limitations as a material representation~\cite{SchwerdtfegerWales2024}.
This work assembles the retained nonlinear lattice energy and conservative
probability equation, then states the conditions needed for cell coupling.
We separate exact conservation identities, conditional analytical results,
reduced-dynamics assumptions and experimental validation requirements.
The objective is a traceable route to crack-initiation probability;
the present framework is not yet a calibrated fatigue-life predictor.

\section{Energy-Based Governing Model}
\subsection{Scope of the reduced model}
Our immediate objective is a probability of first local opening, followed by
an independently validated connection to experimentally observed crack initiation.
Those events must be distinguished throughout the construction.

Let $q^*=(a^*,s^*)$ describe normal separation and unwrapped registry displacement.
Stars denote the original dimensionless variables; $\tau$ is model time.
For one upper coordinate interacting with a fixed infinite lower row,
\begin{align}
 S_p(a^*,s^*)&=\sum_{j\in\mathbb Z}
 \left[(a^*)^2+((j+\tfrac12)b^*-s^*)^2\right]^{-p},\\
 W^*(a^*,s^*)&=4\epsilon^*
 \left[(\sigma_{\rm LJ}^*)^{12}S_6-(\sigma_{\rm LJ}^*)^6S_3\right].
\end{align}
Here $\sigma_{\rm LJ}^*$ is a pair-potential length parameter, not a stress.
The Poisson--Bessel representation evaluates this same infinite sum;
it does not introduce a new interaction law or remove its nonzero lattice modes.
With $a_0^*$ the declared reference separation, define
\begin{equation}
 h^*=(a^*-a_0^*)+\chi s^*,\qquad
 e(q^*)=h^*/a_0^*,\qquad G^*=W^*-f^*(\tau)h^*.
\end{equation}
The observable $e$ is the original reduced total axial strain.
Adding the same elastic strain $\sigma_x/E_x$ to it again would double count
the response already represented by this definition.

\subsection{Probability conservation and first opening}
The retained stochastic hypothesis is overdamped Markov diffusion with constant
positive mobilities and Brownian thermal forcing. Its density satisfies
\begin{equation}
 \partial_\tau P=-\nabla_*\!\cdot J,\qquad
 J=-M^*\left(P\nabla_*G^*+\theta^*\nabla_*P\right),\quad
 M^*=\operatorname{diag}(M_a^*,M_s^*),\quad\theta^*>0.
 \label{eq:retained-fpe}
\end{equation}
Here $P$ is a raw density with respect to the native state measure $dq^*$,
not per unit physical material volume. Under a constant coordinate map
$q_{\rm phys}=Dq^*$, the density at corresponding times transforms as
$P_{\rm phys}(q_{\rm phys})=P(D^{-1}q_{\rm phys})/\det D$.
For $N$ site blocks the joint Jacobian is $(L_aL_s)^N$;
one-site and pair marginals instead use their own two- and four-coordinate measures.
This is an assumed reduced dynamics. Neglect of inertia, loss of memory and
Brownian noise have not yet been justified for these same local coordinates
by an accepted atomistic limit or kinetic calibration. Projection-based
statistical mechanics gives a route through memory-dependent reduced
equations~\cite{Mori1965}; it does not by itself justify the Markov or
overdamped limits for the present coordinates.
The nonlinear energy is retained; a Gaussian approximation is not required
to state or evolve Eq.~\eqref{eq:retained-fpe}.
The initial ensemble must retain its declared source and conditioning.
An equilibrium-basin initialization is a model choice, not a measured defect distribution.

\paragraph{Continuous background and numerical boundaries.}
The positive-$a^*$ energy kernel and unwrapped $s^*$ specify a coordinate
input domain, not a complete continuous probability-background contract.
In the inspected source, the finite SG box uses numerical outer no-flux
boundaries; it does not identify the intended continuous lower/outer
boundaries, reflection field or regulator.
The inspected fixed-interval reductions and historical boundary alternatives
do not supply that contract for the retained full two-coordinate law.
Local smoothness of the exact energy and finite-grid normalization do not
certify global existence, non-explosion or uniqueness of that law.
The continuum initial-law binding and global agreement of the numerical
opening mask with the exact event remain separate requirements.
These are unresolved source bindings, not a proof of model failure.
No numerical cutoff or localization exit is adopted as a new physical wall
or opening event.

Let $D(\tau)$ be the intact domain of the original opening predicate, with
opening boundary $\Gamma(\tau)$. The first exit time from that predicate is
$\tau_{\rm o}$. With a unit initial probability ledger, including any initial
absorbed mass, the raw surviving density gives
\begin{equation}
 S(\tau)=\int_{D(\tau)}P\,dq^*,\qquad
 F_{\rm o}(\tau)=\Pr(\tau_{\rm o}\leq\tau)=A(\tau)=1-S(\tau).
\end{equation}
Absorbed probability is recorded irreversibly and is not normalized back into $P$.
On a regular moving boundary, with outward normal $n$ and velocity $v_\Gamma$,
\begin{equation}
 \dot A=\int_\Gamma (J-Pv_\Gamma)\cdot n\,d\Sigma,\qquad \dot S=-\dot A.
\end{equation}
This smooth-boundary identity does not replace the original discrete rule for
mask changes or instantaneous removal. Reflecting boundaries carry no outward
loss. Numerical repair or remapping loss must not be relabelled physical opening.
The definition of $F_{\rm o}$ alone establishes neither a finite mean opening
time nor equality with an optically detected crack-initiation probability.

\subsection{From model time to cycles and observed initiation}
If a physical clock $t=t_0\tau$ has been verified, let $t$ denote elapsed
physical time from the declared model initialization. Let $\nu(t)$ be the
synchronized recorded frequency in cycles per second (Hz), and $N_{\rm init}$
the recorded cycle index at that origin. Define the cycle count by
\begin{equation}
 \mathcal N(t)=N_{\rm init}+\int_0^t\nu(u)\,du,\qquad
 F_{{\rm o},N}(n)=F_{\rm o}\!\left(t(n)/t_0\right),
\end{equation}
when the measured cycle map has the required inverse $t(n)$.
The cycle count $\mathcal N$ is distinct from the site count $N$.
The experimental time and cycle offsets must be supplied; they are not
silently reset to zero. For constant recorded frequency,
$\mathcal N=N_{\rm init}+\nu t_0\tau$, with physical period $1/\nu$
and model period $1/(\nu t_0)$. Neither $t_0$, $N_{\rm init}$ nor an
unrecorded frequency is estimated here.
Even with this clock, $F_{{\rm o},N}$ becomes an observed crack-initiation CDF
only after the original opening event has been linked to the specified
experimental detection event. This event validation is separate from time conversion.

\subsection{Correlated regions and survivor selection}
For the original $N$-site ($2N$-coordinate) extension, upper--upper interactions remain in
the joint energy in addition to $\sum_iW^*(q_i^*)$; each pair is counted once.
They cannot be removed by dividing space into numerical cells.
Let $P_N$ be the joint density killed when any declared local opening occurs.
Its tagged marginals $u_i$ and $u_{ij}$ both integrate to the same joint survival
$S$, whereas $p_i=u_i/S$ integrates to one when $S>0$.
For the stated block mobility and product intact domain, under the required
regularity and boundary trace assumptions, the exact tagged balance is
\begin{equation}
 \partial_\tau u_i=-\nabla_i\!\cdot j_i-k_i,
\end{equation}
where $k_i$ represents removal when another coordinate opens. The current $j_i$
contains the pair marginal $u_{ij}$ weighted by the pair-energy gradient.
Thus a one-site (two-coordinate) marginal is not a closed independent-cell equation.
The full joint opening-boundary flux already includes all coordinates;
adding the tagged sink $k_i$ to that full flux would count the same loss twice.
Products of conditional marginals are additional closures, not exact joint laws.
No independence-based specimen probability is adopted here without its own
geometrical, dependence and event validation.

\section{Numerical Implementation}
\subsection{Finite volumes in probability space}
The existing two-coordinate reference uses conservative finite volumes and
Scharfetter--Gummel exponential fitting~\cite{ScharfetterGummel1969}.
For neighbouring density values $P_L,P_R$ separated by $\Delta q_\alpha^*$,
the retained interior face formula is
\begin{equation}
 J_{LR}=\frac{M_\alpha^*\theta^*}{\Delta q_\alpha^*}
 \left[B(z)P_L-B(-z)P_R\right],\qquad
 z=\frac{G_R^*-G_L^*}{\theta^*},\quad B(z)=\frac{z}{e^z-1},\quad B(0)=1.
\end{equation}
Opposite face contributions cancel between neighbouring bins. On a fixed
landscape the zero-current ratio is the discrete Gibbs ratio.
Positivity also requires the appropriate time-integration conditions;
conservation alone does not certify a correct distribution or opening CDF.
In particular, bin mass and fully integrated interior-face transfer are
\begin{equation}
 p_K=\int_K P\,dq^*,\qquad
 I_f=\int_{\tau_n}^{\tau_{n+1}}\int_f J\cdot n_f\,d\Sigma\,d\tau.
\end{equation}
The pointwise current $J_{LR}$ above is distinct from $I_f$.
Bin probability mass already includes its quadrature weight. Multiplying it
again by a state-space bin volume or a physical cell volume is incorrect.
An integrated face transfer is probability mass, not a rate to be multiplied
by the time step a second time.

\subsection{Interface to spatial solid mechanics}
Spatial cells and probability bins represent different partitions.
For an explicitly supplied displacement reconstruction, let
$\epsilon_x(U)=B_xU+g_x$ be the spatial axial strain. If $U$ is physical
displacement in metres, $B_x$ has units of inverse metres, while $g_x$ and
strain are dimensionless.
Equating it to the survivor mean is a separate constitutive closure,
\begin{equation}
 R_{\rm strain}=\epsilon_x(U)-\bar e,\qquad
 \bar e=\frac{\int_D eP\,dq^*}{S},\quad S>0.
\end{equation}
The required geometry, reconstruction, face tractions and other stress components
are not determined by the two-coordinate energy. In particular, one axial
response does not supply a general three-dimensional stress law.
For a fixed observable and only the opening sink, its exact survivor mean obeys
\begin{equation}
 \dot{\bar e}=\frac1S\int_D\nabla_*e\cdot J\,dq^*
 -\frac1S\int_\Gamma(e-\bar e)(J-Pv_\Gamma)\cdot n\,d\Sigma.
\end{equation}
The second term is selection of the surviving population. Consequently,
$\dot{\bar e}$ cannot automatically be interpreted as each material cell's
strain rate or mechanical power. Tagged joint marginals additionally require
the corresponding other-opening selection term.
If a coupled trial is repeated, every candidate must start from the same saved
probability state and time interval; only an accepted candidate may advance it.
The current interface does not establish a complete spatial FVM solver or
an atomic commit of displacement, probability and energy histories.

\section{Theoretical Consistency Results}
\subsection{Work and unit consistency}
The existing relaxed axial calibration is $f^*=\kappa\sigma_x/E_x$.
Write $G_{\rm ref}^*$ for the unchanged original single-representative tilted
energy, and let $E_0$ be its energy scale in joules per representative.
For $r$ coherent contributions sharing the same coordinate, the declared
physical total energy would be $G_{\rm coh}=rE_0G_{\rm ref}^*$.
Its virtual work gives
\begin{equation}
 \delta W_{\rm micro}=rE_0f^*a_0^*\delta e,\qquad
 V_{\rm eff}=\frac{rE_0\kappa a_0^*}{E_x}.
\end{equation}
Equality with $V_{\rm geom}\sigma_x\delta e$ requires
$V_{\rm geom}=V_{\rm eff}$ under this mapping. This is a conditional work
normalization, not a measurement of the physical cell volume or atom count.
The coefficient conjugate to dimensionless $q^*$ has units of energy;
for a common length scale $q=L_0q^*$, division by $L_0$ gives its force in newtons.
The factor $r$ is neither the number of independent failure sites nor
automatically the coordinate count $N$.

\paragraph{Coordinate and work-basis binding.}
For a supplied admissible variation map, $\delta x=K_{xq}\delta q$ and
$Q_q=K_{xq}^{T}F_x$ must hold on the same work basis.
The model gradient order $(a^*,s^*)$ and native Cartesian component order
$(s,a)$ must be explicitly distinguished. The coordinate map $K_{xq}$ is
neither the face-jump operator $J_C$ nor supplied by the cell/face ID and
sign relabellings $P_{\rm cell},S_{\rm face}$.
Source energy, coordinate units, counting and representative basis must be
shared by the gradient and every work slot, with each scale applied once.
Environment-translation and environment-residual work are nonadditive views
of the already counted $\sum_i W_i$ contribution; mobile--mobile residual work
is a view of the already counted mobile--mobile pair-energy contribution.
Any separate moment or rotation view must replace a declared disjoint part
of its corresponding residual work, with matching source-component and work-sign tags.
These decompositions are never added to the already included total-gradient work.
The actual map and scales remain unspecified; SI force and three-dimensional
traction therefore remain unestablished without the required physical
length, embedding and face-area data.

Three operations must be distinguished: a passive change of units, a
per-contribution normalization of an already specified total equation, and
a physical change in coherent contribution count.
Dividing an existing equation by its count requires consistent transformations
of its energy, thermal energy, mobility/noise and initial probability measure.
Changing the physical count $r$ at fixed physical temperature is a change
in the energy-to-thermal-energy ratio; it is not a passive unit conversion.
Equality of virtual work or dimensions alone therefore does not establish
the same probability generator or first-opening CDF. Any such equivalence
also requires an explicit representation binding for the energy, thermal
energy, mobility/noise, initial law and opening rule. No missing scale or
normalization mode is selected by the work identity.

The following passive formulas are restricted to the unchanged $r=1$
reference representation, so that $G_{\rm phys}=E_0G_{\rm ref}^*$.
With constant $D=\operatorname{diag}(L_a,L_s)$ and positive constant
$E_0,t_0$, let $q_{\rm phys}=Dq^*$ and $t_{\rm phys}=t_0\tau$.
For the original reference mobility $M^*$ the same generator requires
\begin{equation}
 k_BT_{\rm phys}=E_0\theta^*,\qquad
 M_{\rm phys}=\frac{DM^*D^T}{t_0E_0}.
\end{equation}
Here $M_{\rm phys}$ has units $\mathrm{m^2/(J\,s)}$.
The initial measure, boundary and density Jacobian must be transported as well.
Both coordinate mobilities must yield the same $t_0$ for this same energy
representation. No accepted local mobility calibration presently supplies
that clock; model time is therefore not labelled seconds or experimental cycles.

\paragraph{Clock compatibility and identifiability.}
Shared-clock compatibility does not itself identify physical seconds.
For conditionally fixed $E_0,D$ and a compatible candidate
$(t_0,M_{\rm phys})$, any dimensionless $c>0$ gives
\begin{equation}
\begin{aligned}
 t_0^{(c)}&=ct_0,\\
 M_{\rm phys}^{(c)}&=M_{\rm phys}/c .
\end{aligned}
\end{equation}
The native law, initial ensemble, opening rule and $f^*(\tau)$ remain
fixed. These are distinct, uncalibrated physical embeddings, not a
passive identity or verified dynamics of the same material or experiment.
The physical-time placement of the load changes with $c$. The native
first-opening CDF is unchanged, but equality of the corresponding
first-opening CDFs at a fixed physical time is not implied.

With independently established $E_0,D$ and absolute physical mobility,
a positive diagonal component ratio in the native-coordinate pullback
fixes a candidate clock for those inputs.
The full matrix relation, including cross terms, must also hold for compatibility;
agreement of the two axis clocks without those independent inputs
does not identify physical seconds.
Temperature and static anchors do not fix the kinetic clock.
Same-test frequency $\nu$ and an independently supplied and verified native
full period $P_\tau$, with the load waveform, cycle count, phase, and
initialization/time origins matched to that same test, can supply the
protocol-map candidate $t_{0,\rm protocol}=(1/\nu)/P_\tau$.
That map does not independently validate the mobility or generator timing.
The actual anchors remain unspecified; a single stochastic onset is
not equated deterministically to one native first-passage time.

\subsection{A conditional interior probability law}
Here $E=\{(a^*,s^*):a^*>0,\ s^*\in\mathbb R\}$ denotes the
energy-input domain only. On a predeclared finite native interval $[0,T]$,
let the deterministic load satisfy $f^*\in L^1(0,T)$ and retain the original
positive finite constant mobilities and temperature parameter.
Zero denotes the existing initialization reference, not a newly chosen
experimental origin. For any open $U$ with compact closure inside $E$, set
\begin{equation}
 \sigma_U=\inf\{\tau\in[0,T]:Q_\tau\notin U\},\qquad
 \inf\varnothing=\infty .
\end{equation}
Consider the exact native interior equation
\begin{equation}
\begin{aligned}
 dQ_\tau&=\left[-M^*\nabla_*W^*(Q_\tau)
       +M^*(1,\chi)^T f^*(\tau)\right]d\tau+\Sigma^*dB_\tau,\\
 \Sigma^*&=(2\theta^*M^*)^{1/2}.
\end{aligned}
\end{equation}
For the same predeclared Borel probability initial law $\mu$ on $E$,
assume an initial variable measurable at initialization and a Brownian
motion for the same usual filtration. There is a unique continuous law
of paths stopped at $\sigma_U$ among the stated interior candidates,
with no regulator, additional drift or sink inside $U$.
It is independent of the auxiliary cutoff used in its construction and
admits a Borel initial-to-stopped-path kernel $K_U$.
An initial point outside $U$ gives $\sigma_U=0$ and the constant path at
that same point; no initial mass is discarded or reconditioned.

To see this, reuse local smoothness of the exact energy and choose a smooth
spatial cutoff equal to one near $\overline U$, supported compactly inside
$E$. Its zero extension gives a bounded globally Lipschitz auxiliary
spatial drift. The $L^1$ load has an absolutely continuous primitive;
pathwise Picard iteration gives a continuous nonanticipative solution map
for the additive driving path. Exit from open $U$ hits its closed
complement when finite on this horizon. A distance-minimum argument gives
a Borel stopping map without assuming a continuous hitting-time functional
or a regular boundary. Gronwall comparison until the minimum of two exits,
followed by exit attainment, proves cutoff independence and stopped
pathwise uniqueness. The same construction on a weak candidate's space
identifies its law, and composition with the Borel stopping map gives $K_U$.

Let $\mathcal C_T=C([0,T],\mathbb R^2)$. The kernel mixture extends
linearly to the same predeclared finite raw initial measure $\mu$:
\begin{equation}
\begin{aligned}
 \Pi_\mu(A)&=\int_E K_U(x,A)\,\mu(dx),\\
 \Pi_\mu(\mathcal C_T)&=\mu(E),\qquad
 \Pi_\mu\{\sigma_U=0\}=\mu(E\setminus U).
\end{aligned}
\end{equation}
For a raw measure, $\Pi_\mu$ is a path measure and is not forced to have
unit mass. Any separately recorded initial absorbed mass retains its
original ledger. Localization does not add to physical absorption.
This is a conditional interior theorem, not certification that a supplied
original reflected or killed law has this restriction. That application
still needs its original background, initial-law and regulator bindings.
The auxiliary continuation after an earlier original opening is not
validated by this construction. No actual $U$, wall, load or physical
scale is selected; global non-explosion, $a^*=0$ accessibility, full
reflection and moving/table opening-event equivalence remain outside
the result.


\subsection{Localization and a conditional opening-mass ledger}
Retain the same finite raw initial measure $\mu$, native origin and
stopped path measure $\Pi_\mu$, with $m=\mu(E)$. Localization freezes
the coordinate after $\sigma_U$; applying a later load-dependent mask to
that frozen point can produce a signal without identifying the original
post-exit motion. Such a signal is not counted as original opening.
Moreover, $f^*\in L^1$ specifies the interior drift through its time
integral, whereas an opening predicate may depend on pointwise load,
jump-side, contact, branch or table conventions. The original pointwise
load and event rules must therefore be supplied separately. Neither
continuous coordinates nor the interior theorem certifies an arbitrary
moving or tabulated opening event.

For a finite native observation time $t\in[0,T]$, suppose the original
rule supplies measurable certified prefix events $O_t$ and $V_t$.
In $O_t$, first opening is verified at $u<t$ and $u<\sigma_U$;
in $V_t$, $t<\sigma_U$ and absence of opening is verified through and
including $t$. An unresolved initial or local predicate is in neither
class. These events are disjoint. With $C_t=(O_t\cup V_t)^c$ and
$o_t=\Pi_\mu(O_t)$, $v_t=\Pi_\mu(V_t)$, $c_t=\Pi_\mu(C_t)$,
finite-measure additivity gives
\begin{equation}
 m=o_t+v_t+c_t.\label{eq:localization-mass-ledger}
\end{equation}
These are integrated masses on the original measure, without survivor
normalization or extra state, cell-volume or time weights.
If a measurable source certificate $G_t\subseteq\{Z_0\in U\}$ is
supplied with $O_t,V_t\subseteq G_t$, then
\begin{equation}
\begin{aligned}
 C_t&=X_0\mathbin{\dot\cup}J_t\mathbin{\dot\cup}C_{G,t},\\
 X_0&=\{Z_0\notin U\},\quad J_t=\{Z_0\in U\}\setminus G_t,\\
 C_{G,t}&=G_t\setminus(O_t\cup V_t),\\
 c_t&=\mu(E\setminus U)+\Pi_\mu(J_t)+\Pi_\mu(C_{G,t}).
\end{aligned}
\end{equation}
An absent certificate leaves actual $O_t,V_t$ and their masses
unresolved, rather than zero. Initial outside-$U$ mass and other
unobserved mass are not added to physical absorption.

A finer race partition requires source-supplied original joint stopping
records, or the same decidable first-event race, on $G_t$. Only under
that additional contract, writing $\tau_o$ for original first opening
and $\sigma=\sigma_U$, one has the disjoint partition
\begin{equation}
\begin{aligned}
 O_t&=G_t\cap\{\tau_o<\sigma,\ \tau_o<t\},\\
 L_t&=G_t\cap\{\sigma<\tau_o,\ \sigma<t\},\\
 B_t&=G_t\cap\{\tau_o=\sigma<t\},\\
 H_t&=G_t\cap\{\min(\tau_o,\sigma)=t\},\\
 V_t&=G_t\cap\{t<\tau_o,\ t<\sigma\}.
\end{aligned}
\end{equation}
Indeed, the minimum is below, equal to, or above $t$; below $t$ the
two records have either strict order or a tie. Thus
$C_{G,t}=L_t\mathbin{\dot\cup}B_t\mathbin{\dot\cup}H_t$.
Simultaneous and observation-endpoint ties retain their mass; no new
priority or zero-tie assumption is made. For finite $t$, two infinite
records belong to $V_t$. A killed coordinate without the required joint
record does not justify inventing continuation after opening; only the
coarser certified ledger then remains available.

To connect this ledger to an original global raw law $\mathbb Q_\mu$,
the source must separately identify the same prefix events with masses
$o_t,v_t$, preserve total mass $m$, and certify
$O_t^{\rm glob}\subseteq\{\tau_o\leq t\}$ and
$V_t^{\rm glob}\subseteq\{\tau_o>t\}$. The same initial measure alone
does not establish that restriction. For their global complement
$C_t^{\rm glob}$, of mass $c_t$, and a separately supplied initial
absorbed cohort $A_0$, the conditional identity and bounds are
\begin{equation}
\begin{aligned}
 A_{\rm raw}(t)&=A_0+o_t+
 \mathbb Q_\mu(C_t^{\rm glob}\cap\{\tau_o\leq t\}),\\
 A_0+o_t&\leq A_{\rm raw}(t)\leq A_0+o_t+c_t=A_0+m-v_t.
\end{aligned}
\label{eq:conditional-opening-interval}
\end{equation}
Here $A_0$ and $\mu$ must share the original preparation, mass basis,
origin and event meaning and have disjoint cohorts, including consistency
of any $\tau_o=0$ mass. Together with those source bindings, the original
probability ledger $A_0+m=1$ is required before interpreting
$A_{\rm raw}$ as the original first-opening CDF. Otherwise the result
remains a raw-mass interval and is not arbitrarily divided to make a
probability. The residual $c_t$ is not a supplied numerical error or a certified
boundary-error estimate. No claim is made that the conditional interval
is sharp or that every unobserved completion is attainable.
The interval's upper endpoint $A_0+o_t+c_t$ is not selected as an actual CDF value.
Actual certificates, masses, original global restriction and physical
clock remain unspecified.


\paragraph{Initial mass in numerical output.}
At the literal local snapshot \texttt{5f87837}, the numerical source defines
\texttt{cumulative\_\allowbreak absorbed\_\allowbreak mass} and
\texttt{initiation\_\allowbreak probability} to include initial removal.
Treating either as $o_t$ and adding the same $A_0$ would double-count that cohort.
Repeated snapshots of \texttt{initial\_\allowbreak absorbed\_\allowbreak mass}
record an initial ledger value, not new increments.
Post-initial integrated flux is a different aggregate; excluding initial
removal does not certify it as strict-interior original opening mass $o_t$.
Density, rate and integrated mass retain distinct measures and intervals;
quadrature, state-Jacobian, physical-volume and time factors must not be applied twice.
Consumption requires source bindings for the preparation and alive raw
initial-measure stage, normalization history, matching disjoint cohorts
including zero-time event accounting,
initialization origin, event meaning and eligible prefix.
With $m=\mu(E)$, the original probability ledger $A_0+m=1$ and those
bindings plus the global prefix restriction are required before interpreting
the conditional raw-mass interval as a first-opening CDF;
normalization alone is insufficient.
Without those bindings, the result remains a raw-mass interval without renormalization.
Numerical residual fields do not supply the unresolved mass $c_t$.
Source definitions alone do not establish an actual run or its raw initial population.
Actual raw $\mu$, execution records, cohort bindings and the physical clock remain unspecified.

\paragraph{Time and event records.}
In the inspected local snapshot at commit \texttt{5f87837}, load evaluations are numerical point calls. The explicit force at \(t\) and the implicit force at \(t+\Delta t\) identify integration stages; they do not select left or right limits at a load jump. Energy, mask, subsequent loop, and final record evaluations make separate calls, even when their time arguments coincide. Without a supplied original pointwise load representative and a repeatability contract, these calls are not certified evaluations of the same original \(f(t)\). Numerical record and duration tolerances, together with the last snapshot, resolve neither \(\tau_{\rm o}<t\) versus \(\tau_{\rm o}=t\) nor the opening, localization, and simultaneous cases at the observation horizon; they provide no CDF error bound. Aggregate probability, removal, and flux fields do not supply a joint \((\tau_{\rm o},\sigma_U)\) event order, the certificate \(G_t\), or a restriction identifying the same first-event prefix of an original global law. These observations concern the inspected source and interface and the records available to this review. They establish no actual-run validation or SG--continuum event equivalence, and imply neither repository-wide absence of records nor model failure. The original load and opening rules remain unchanged.

\subsection{Conditional finite-time load sensitivity}
Here $T$ and the integration variable $u$ denote the original model time,
with no conversion to physical seconds.
Consider two prescribed deterministic loads $f_0^*,f_1^*$ on $[0,T]$,
with the same $W^*$, positive constant $M^*,\theta^*$, initial law,
model clock and background reflection rule. Assume existence, non-explosion,
the weak uniqueness needed to identify the changed law, and a common measurable
opening event. If $\Delta f^*=f_1^*-f_0^*\in L^2(0,T)$, the previously
derived finite-time comparison is
\begin{equation}
 |F_1^D(T)-F_0^D(T)|\leq B_{\rm load}(T)
 :=\min\left\{1,\sqrt{\frac{M_a^*+\chi^2M_s^*}{8\theta^*}
 \int_0^T |\Delta f^*(u)|^2\,du}\right\}.
\end{equation}
The drift difference is $M^*(1,\chi)^T\Delta f^*$. A true-martingale
change of measure on the driving space gives an entropy upper bound;
data processing and Pinsker's inequality then bound any common path event.
For deterministic square-integrable drift control the Novikov exponential
is finite, but this does not establish the assumed global well-posedness of
the singular lattice/reflection problem. Entropy results require their
stated hypotheses~\cite{Leonard2012,Lehec2025}; the upper bound is not
obtained by reversing an entropy inequality.
When the initial law or load-dependent opening predicate differs, the
comparison instead requires
\begin{equation}
 |F_1(T)-F_0(T)|\leq
 \min\{1,\eta_{\rm init}+B_{\rm load}(T)+\eta_{\rm event}(T)\},
\end{equation}
where $\eta_{\rm init}$ is the initial-law total variation distance and
$\eta_{\rm event}$ is the probability of the symmetric difference of the
two opening events under the chosen common comparison law.
The load integral alone does not control the moving-boundary term.
This is an absolute CDF bound for finite time; it is not a relative rare-event
error, an infinite-time mean-lifetime bound or a physical calibration.
No numerical value for its inputs or event term is supplied.

\subsection{A conditional exact far-field result}
In this paragraph stars are suppressed in the original lattice units.
For positive $\epsilon,\sigma_{\rm LJ},b,a$, define the zero mode
$W_0=63\pi\epsilon\sigma_{\rm LJ}^{12}/(64ba^{11})
-3\pi\epsilon\sigma_{\rm LJ}^{6}/(2ba^5)$.
The reviewed exact-sum bounds, uniform for every real $s$, are
\begin{align}
 |S_p-A_p|&\leq a^{-2p},\quad
 A_p=\frac{\sqrt\pi\,\Gamma(p-1/2)}{b\Gamma(p)}a^{1-2p},\quad p>\tfrac12,\\
 |W_a-W_0'|&\leq24\epsilon\sigma_{\rm LJ}^6
 (a^{-7}+2\sigma_{\rm LJ}^6a^{-13}),\\
 |W_{aa}-W_0''|&\leq24\epsilon\sigma_{\rm LJ}^6
 (9a^{-8}+30\sigma_{\rm LJ}^6a^{-14}).
\end{align}
The first bound follows by partitioning the line into intervals centred at
the shifted lattice sites and bounding the quadrature error by half the total
variation of $(a^2+x^2)^{-p}$. Exact derivatives of the lattice sum, followed
by separate remainder bounds, give the other two inequalities; an error
inequality is not differentiated. In particular,
\begin{equation}
 \frac{2079\pi}{16}\left(\frac{\sigma_{\rm LJ}}a\right)^6
 +216\frac ba+720\frac ba\left(\frac{\sigma_{\rm LJ}}a\right)^6<45\pi
 \quad\Longrightarrow\quad W_a>0,\quad W_{aa}<0.
\end{equation}
This sufficient condition gives at most one fixed-$s$ scalar root of $W_a=f$
within the declared tail. It does not establish root existence, a full
two-dimensional saddle, the operating selector, or a numerical boundary table.
No actual threshold or evaluated finite-tube curvature bound has been
certified for an operating parameter/domain instance. The conditional
symbolic finite-tube result below is separate from such instance certification.
Indeed $W_{aa}\to0$ in the far field, so strict negative curvature there
does not provide a positive uniform curvature bound on the entire infinite tail.

\subsection{A conditional finite root tube}
Continue in the same native lattice units, with stars suppressed.
Let $0<R<A<\infty$ and $0<f_{\min}\leq f_{\max}<\infty$.
The constant $A$ here is the upper root-tube endpoint, not the absorbed
probability $A(\tau)$. Reusing the exact far-field envelopes, define
\begin{align}
 F_0(a)&=W_0'(a)
 =\frac{15\pi\epsilon\sigma_{\rm LJ}^6}{2b}a^{-6}
 -\frac{693\pi\epsilon\sigma_{\rm LJ}^{12}}{64b}a^{-12},\\
 E_F(a)&=24\epsilon\sigma_{\rm LJ}^6
 (a^{-7}+2\sigma_{\rm LJ}^6a^{-13}),\\
 L_F(a)&=F_0(a)-E_F(a),\qquad U_F(a)=F_0(a)+E_F(a),\\
 M(a)&=45\pi-\frac{2079\pi}{16}\left(\frac{\sigma_{\rm LJ}}a\right)^6
 -216\frac ba-720\frac ba\left(\frac{\sigma_{\rm LJ}}a\right)^6 .
\end{align}
Here $M(a)$ is a scalar curvature margin, distinct from mobility $M^*$.
A reviewed sufficient set of symbolic conditions is
\begin{equation}
 M(R)>0,\qquad L_F(R)>f_{\max},\qquad U_F(A)<f_{\min}.
 \label{eq:finite-tube-conditions}
\end{equation}
Define $\eta_L=L_F(R)-f_{\max}>0$ and
$\eta_R=f_{\min}-U_F(A)>0$. The endpoint inequalities are then
$W_a(R,s)\geq f_{\max}+\eta_L$ and
$W_a(A,s)\leq f_{\min}-\eta_R$ for every real $s$.
On the whole finite tube the corresponding conservative symbolic lower bound is
\begin{equation}
 -W_{aa}(a,s)\geq k_{\rm cert}^{\rm sym}(R,A)
 :=\frac{\epsilon\sigma_{\rm LJ}^6}{bA^7}M(R)>0,
 \qquad R\leq a\leq A .
 \label{eq:finite-tube-curvature}
\end{equation}
These conditions supply the strict bracket and curvature inputs to the
existing scalar root-graph theorem for a conditional root
$a_{\rm tube}(s,f)\in(R,A)$, $f\in[f_{\min},f_{\max}]$.
They are sufficient conditions, not assertions that every operating band
or arbitrary choice of endpoints satisfies them.
For fixed admissible $R$, $k_{\rm cert}^{\rm sym}(R,A)\to0$ as
$A\to\infty$, consistently with the absence of a positive uniform bound
on the entire infinite tail.

For two force histories confined to this band and roots on this same tube
graph, the existing same-slip boundary comparison can use
\begin{equation}
 |\Delta a_{\rm tube}(\tau,s)|
 \leq\frac{|\Delta f(\tau)|}{k_{\rm cert}^{\rm sym}}
 \leq\frac{\delta}{k_{\rm cert}^{\rm sym}},
 \qquad |\Delta f(\tau)|\leq\delta .
\end{equation}
Its initial-law, background, reflection and event hypotheses remain required.
Identifying this tube graph with the original mechanical outer boundary
requires separate selector and left-well evidence; otherwise the estimate
concerns the tube graph alone. It is a native boundary-location bound,
not a distance between stochastic paths, a CDF bound, a measured crack
length or an experimental error bound.
The tube introduces no probability-domain truncation, reflecting wall or
change of initial ensemble. Actual $R,A$, force band, slip domain,
parameter/runtime binding and evaluated $k_{\rm cert}$ remain unspecified;
no actual branch-certificate instance has passed.
The full two-dimensional saddle, operating selector and left well,
finite-evaluator/roundoff/table/SG accuracy and physical calibration remain
uncertified.

\paragraph{Symbolic nonvacuity and order of declaration.}
The finite-tube conditions admit a nonempty symbolic family.
For fixed positive $\epsilon,\sigma_{\rm LJ},b$, put
$x=b/a$, $y=(\sigma_{\rm LJ}/a)^6$,
$C(a)=\epsilon\sigma_{\rm LJ}^6/(ba^6)>0$ and
$\ell(a)=L_F(a)/C(a)$. The reviewed exact identity is
\begin{equation}
 \ell(a)=\frac{M(a)}6+\frac{693\pi}{64}y+12x+72xy .
\end{equation}
Thus $M(R)>0$ implies $L_F(R)>0$.
Since $M(a)\to45\pi$ as $a\to\infty$, admissible finite $R$ exist.
After an admissible $R$ is declared, declare a band
$0<f_{\min}\leq f_{\max}<L_F(R)$.
For each such positive band, finite upper endpoints exist; it is sufficient to require
\begin{equation}
\begin{aligned}
 A&>\mathcal A(R,f_{\min}),\\
 \mathcal A(R,f_{\min})&=\max\left\{
 \begin{gathered}
 R,\quad
 \left(\frac{45\pi\epsilon\sigma_{\rm LJ}^6}{2bf_{\min}}\right)^{1/6},\\
 \left(\frac{72\epsilon\sigma_{\rm LJ}^6}{f_{\min}}\right)^{1/7},
 \quad
 \left(\frac{144\epsilon\sigma_{\rm LJ}^{12}}{f_{\min}}\right)^{1/13}
 \end{gathered}\right\}.
\end{aligned}
\end{equation}
Dropping the negative $a^{-12}$ term from $U_F$ gives a positive
majorant; this strict threshold bounds each of its three positive
terms by less than $f_{\min}/3$, and hence $U_F(A)<f_{\min}$.
The order is positive parameters, an admissible $R$, a band below
its positive cap, and then a finite $A$ depending on that band.
This neither certifies an arbitrary preselected operating band nor
supplies one finite $A$ for all bands approaching zero force.
Zero or negative force, the actual selector and left well,
finite-evaluator/CDF accuracy and physical crack initiation remain
uncertified. No actual parameters or endpoints are selected, and the
probability domain, initial law and opening event are unchanged.

\section{Discussion}
Earlier one-dimensional and Weibull-based constructions remain in the archived
draft. They are not identified with the present two-coordinate dynamics, and
their illustrative fitted values are not adopted as current-model inputs.
The same pair energy and the same external work must each be counted once.
An API returning $G$ already includes the loading tilt; subtracting $fh$ again
would change the model. Likewise the macro and micro virtual works above
describe the same transfer, rather than two additive stored energies.
The lattice zero mode and registry modes are a mathematical decomposition;
they have not been identified with physical bulk and surface energies.
Such an excess decomposition requires the same atom set, energy reference,
constraints, strain path and boundary conditions on both sides of the subtraction.
For the retained reduced model that matched physical bulk reference is unresolved.
Ideal uniform-interface strength, defect-controlled yielding and first crack
initiation consequently remain different observables.

Experimental validation requires records from the same specimen and test:
actual cycle indices and load frequency, continuous load or a verified
synchronization map, the last negative and first positive observation,
and the detection criterion, correction and resolution.
Those observations can bracket a detected transition under the stated inspection
assumptions. They do not by themselves establish when a sub-resolution crack
first formed or quantify the detection delay. The same-test controller, calibration and negative--positive records
needed for the proposed validation are not yet available to this study.
No material lifetime or observed initiation probability is inferred here.

\section{Conclusion}
The current result is a conservative first-opening probability framework with
explicit energy, population and work definitions. Its nonlinear reference
energy and opening rule are preserved. The exact far-field result supplies
a conditional analytical constraint, while the physical clock, matched bulk
reference, spatial constitutive closure and observed-crack mapping remain
separate validation tasks. These missing inputs are left unspecified rather
than inferred from illustrative fits or numerical conservation alone.

\bibliographystyle{unsrt}
\bibliography{1}
\end{document}
